\section*{Chapitre \ref{ch:b1:field-lab-methods-1} --- Méthodes de terrain et de laboratoire I}

\begin{solution}{exo:b1:field-lab-methods-1:1}
Au 1:10\,000, \qty{4.9}{m} devient $4.9/10\,000 = \qty{0.49}{mm}$~; au
1:25\,000, $4.9/25\,000 = \qty{0.20}{mm}$, environ la largeur d'un trait de crayon.
\end{solution}

\begin{solution}{exo:b1:field-lab-methods-1:2}
À Reykjavik, $D = \ang{-11.28}$~: $52 - 11.28 = \ang{040.7}$. À Wellington,
$D = \ang{+23.23}$~: $52 + 23.23 = \ang{075.2}$. La même lecture de boussole
signifie des directions écartées de \ang{34.5}.
\end{solution}

\begin{solution}{exo:b1:field-lab-methods-1:3}
$20\to015$. Avec $t = \ang{15}$ et $p = \ang{20}$, le vecteur unitaire descendant est
\[
  (\cos\ang{20}\sin\ang{15},\ \cos\ang{20}\cos\ang{15},\ -\sin\ang{20}) =
  (0.2432,\ 0.9077,\ -0.3420).
\]
\end{solution}

\begin{solution}{exo:b1:field-lab-methods-1:4}
La moyenne arithmétique est $(350 + 358 + 6)/3 = \ang{238}$, presque opposée aux
trois lectures. Les vecteurs unitaires horizontaux $(\sin a, \cos a)$ s'additionnent en
$(-0.1736 - 0.0349 + 0.1045,\ 0.9848 + 0.9994 + 0.9945) = (-0.1040, 2.9787)$,
dont l'\omterm{def:b1:spherical-astronomy:horizontal}{azimut} est $\operatorname{atan2}(-0.1040, 2.9787) = \ang{-2.0}$, c'est-à-dire
\ang{358.0}~: la bonne moyenne, au milieu des lectures.
\end{solution}

\begin{solution}{exo:b1:field-lab-methods-1:5}
Moyenne $175/5 = \ang{35.0}$. Écarts $-1, 2, 0, -2, 1$~; somme des carrés 10~;
$s = \sqrt{10/4} = \ang{1.58}$. Erreur type $1.58/\sqrt5 = \ang{0.71}$.
\end{solution}

\begin{solution}{exo:b1:field-lab-methods-1:6}
Les normales sont $(\sin\ang{30}, 0, \cos\ang{30}) = (0.5, 0, 0.866)$ et
$(-0.5, 0, 0.866)$~; $\vec S = (0, 0, 1.732)$ est vertical~: la moyenne est un
plan horizontal, et $\bar R = 1.732/2 = 0.866$, chaque lecture étant à \ang{30}
de la moyenne. Les deux lectures n'échantillonnent pas une même surface avec des
erreurs~: elles appartiennent à deux flancs d'un \omterm{def:b1:geological-maps:fold}{pli}, deux populations
différentes, et leur moyenne ne décrit ni l'une ni l'autre. Chaque flanc se
moyenne séparément~; ensemble, ils définissent l'\omterm{def:b1:geological-maps:fold}{axe du pli}, la droite commune
aux deux plans, ici horizontale et nord--sud.
\end{solution}

\begin{solution}{exo:b1:field-lab-methods-1:7}
$\tan p = \tan\ang{12}\cos(\ang{80} - \ang{120}) = 0.2126\times0.7660 =
0.1628$~: $p = \ang{9.2}$, les stries plongeant vers \ang{120}. Pour la direction
\ang{170}, $\cos(\ang{-90}) = 0$~: la ligne est horizontale, le long de la
direction de la couche.
\end{solution}

\begin{solution}{exo:b1:field-lab-methods-1:8}
$0.47/0.20 = 2.35$ ans après le 1\textsuperscript{er} janvier 2026, c'est-à-dire
au printemps 2028. Au 1\textsuperscript{er} janvier 2030~:
$-0.47 + 4\times0.20 = \ang{+0.33}$, à l'est du nord géographique.
\end{solution}

\begin{solution}{exo:b1:field-lab-methods-1:9}
Les trois normales sont dans le plan vertical d'\omterm{def:b1:spherical-astronomy:horizontal}{azimut} \ang{080}, à
\ang{10}, \ang{12} et \ang{14} de la verticale~; les deux extrêmes sont
symétriques par rapport à celle du milieu, donc leur somme est le long d'elle~:
la moyenne est $080/12$, avec $\theta = \ang{2}, 0, \ang{2}$. Alors $\bar R = (1 +
2\cos\ang{2})/3 = (1 + 1.998782)/3 = 0.999594$, et $\sqrt{2(1 - \bar R)} =
\sqrt{0.000812} = \qty{0.02850}{rad} = \ang{1.633}$~; exactement,
$\theta_{\text{rms}} = \sqrt{(4 + 0 + 4)/3} = \ang{1.633}$.
\end{solution}

\begin{solution}{exo:b1:field-lab-methods-1:10}
La section est un cercle de rayon $\sqrt{r^2 - h^2}$. Sa moyenne est
$\frac1r\int_0^r\sqrt{r^2 - h^2}\,\mathrm dh = \frac1r\cdot\frac{\pi r^2}{4}
= \frac{\pi r}{4}$, puisque l'intégrale est l'aire d'un quart de disque. Des
grains de \qty{0.50}{mm} de diamètre apparaissent en moyenne avec
$0.50\times\pi/4 = \qty{0.39}{mm}$ de diamètre~: une \omterm{def:b1:field-lab-methods-1:thin-section}{lame mince} sous-estime la
taille des grains.
\end{solution}

\begin{solution}{exo:b1:field-lab-methods-1:11}
La rotation de $D$ autour de la verticale,
$(x, y, z) \mapsto (x\cos D + y\sin D,\ -x\sin D + y\cos D,\ z)$, envoie
$(\sin\delta\sin\alpha, \sin\delta\cos\alpha, \cos\delta)$ sur
$(\sin\delta\sin(\alpha + D), \sin\delta\cos(\alpha + D), \cos\delta)$~: elle
ajoute $D$ à chaque \omterm{def:b1:geological-maps:attitude}{sens du pendage}. Étant linéaire, elle envoie la résultante
$\vec S$ sur la résultante des normales tournées~; elle conserve les longueurs et
la composante verticale, donc le \omterm{def:b1:geological-maps:attitude}{pendage} moyen $\arccos(S_U/\|\vec S\|)$ est
inchangé, et elle fait tourner la partie horizontale de $\vec S$ de $D$, donc le
\omterm{def:b1:geological-maps:attitude}{sens du pendage} moyen augmente de $D$. Corriger de la déclinaison chaque lecture,
ou seulement leur moyenne, donne le même résultat.
\end{solution}

\begin{solution}{exo:b1:field-lab-methods-1:12}
La dispersion le long de la ligne de plus grande pente est une dispersion du
\omterm{def:b1:geological-maps:attitude}{pendage}~: l'erreur type est $2/\sqrt N \le 0.5$ pour $N \ge 16$. En travers, une
rotation de la normale de \ang{2} change le \omterm{def:b1:geological-maps:attitude}{sens du pendage} de
$2/\sin\ang{12} = \ang{9.62}$~; l'erreur type du \omterm{def:b1:geological-maps:attitude}{sens du pendage} moyen est
$9.62/\sqrt N \le 1$ pour $N \ge 92.5$, soit 93 lectures. Un \omterm{def:b1:geological-maps:attitude}{pendage} faible est
vite connu~; son sens ne l'est pas.
\end{solution}

\begin{solution}{pb:b1:field-lab-methods-1:1}
\textbf{1.} \omterm{def:b1:geological-maps:attitude}{Sens du pendage}~: $967/12 = \ang{80.6}$~; \omterm{def:b1:geological-maps:attitude}{pendages}~: $145/12 =
\ang{12.1}$.

\textbf{2.} \ang{355} et \ang{005} ont pour moyenne \ang{180}, la direction
opposée. Ici tous les \omterm{def:b1:geological-maps:attitude}{sens du pendage} sont compris entre \ang{076} et \ang{085},
loin du saut à \ang{000}, donc la moyenne arithmétique convient.

\textbf{3.} $\|\vec n\|^2 = \sin^2\delta(\sin^2\alpha + \cos^2\alpha) +
\cos^2\delta = 1$. La ligne de \omterm{def:b1:geological-maps:attitude}{pendage} descendante est dirigée selon
$(\cos\delta\sin\alpha, \cos\delta\cos\alpha, -\sin\delta)$, et son produit
scalaire avec $\vec n$ est $\sin\delta\cos\delta - \cos\delta\sin\delta = 0$~; la
ligne de direction est horizontale et perpendiculaire au \omterm{def:b1:geological-maps:attitude}{sens du pendage},
$(\cos\alpha, -\sin\alpha, 0)$, et son produit scalaire avec $\vec n$ est
$\sin\delta\sin\alpha\cos\alpha - \sin\delta\cos\alpha\sin\alpha = 0$.

\textbf{4.} La lecture 11 est la même que la lecture 2~: $(0.2222, 0.0352, 0.9744)$.
Lecture 12~:
\[
  (\sin\ang{11}\sin\ang{82},\ \sin\ang{11}\cos\ang{82},\ \cos\ang{11}) =
  (0.1890,\ 0.0266,\ 0.9816).
\]

\textbf{5.} $\vec S = (2.4747,\ 0.4082,\ 11.7294)$.

\textbf{6.} $\|\vec S\| = \sqrt{6.1241 + 0.1666 + 137.5788} = 11.99457$, et
$\bar R = 11.99457/12 = 0.999547$.

\textbf{7.} $\bar\alpha = \operatorname{atan2}(2.4747, 0.4082) = \ang{80.6}$
et $\bar\delta = \arccos(11.7294/11.99457) = \arccos 0.977893 = \ang{12.1}$~:
la moyenne vectorielle est $080.6/12.1$, la même que les moyennes arithmétiques ici.

\textbf{8.} \omterm{def:b1:geological-maps:attitude}{Sens du pendage}~: $s = \ang{2.94}$~; \omterm{def:b1:geological-maps:attitude}{pendages}~: $s = \ang{1.68}$.

\textbf{9.} $\theta_{\text{rms}} \approx \sqrt{2\times0.000453} =
\qty{0.0301}{rad} = \ang{1.72}$.

\textbf{10.} $\sqrt{1.68^2 + (\sin\ang{12.1}\times2.94)^2} = \sqrt{2.82 +
0.38} = \ang{1.79}$, proche de \ang{1.72}~: les lectures se comportent comme le
dit le modèle d'erreur, aux effets près de l'arrondi au degré et de la petitesse
de l'échantillon.

\textbf{11.} Le \omterm{def:b1:geological-maps:attitude}{pendage}~: il contribue pour 2.82 sur 3.20 (en degrés carrés). Le
\omterm{def:b1:geological-maps:attitude}{sens du pendage} ne compte que par $\sin\delta = 0.21$, car sur une couche
faiblement inclinée un changement de \omterm{def:b1:geological-maps:attitude}{sens du pendage} fait très peu tourner le
plan.

\textbf{12.} \omterm{def:b1:geological-maps:attitude}{Pendage} moyen~: $1.68/\sqrt{12} = \ang{0.48}$~; \omterm{def:b1:geological-maps:attitude}{sens du pendage} moyen~:
$2.94/\sqrt{12} = \ang{0.85}$.

\textbf{13.} Non~: \ang{2.9} de \omterm{def:b1:geological-maps:attitude}{sens du pendage} sur une couche de \omterm{def:b1:geological-maps:attitude}{pendage} \ang{12}
ne représentent que $2.9\sin\ang{12.1} = \ang{0.6}$ de rotation du plan, moins
que les \ang{1.7} des \omterm{def:b1:geological-maps:attitude}{pendages}. La moindre ondulation d'une surface de
stratification faiblement inclinée disperse largement son \omterm{def:b1:geological-maps:attitude}{sens du pendage}.

\textbf{14.} $D = \ang{-2.5}$~: $83.1 - 2.5 = \ang{080.6}$~; sa moyenne est
$080.6/12.1$, la même que celle du premier géologue.

\textbf{15.} Le \omterm{def:b1:geological-maps:attitude}{pendage} se lit avec le clinomètre, par rapport à la verticale
donnée par la pesanteur~; l'aiguille n'y joue aucun rôle.

\textbf{16.} Oui~: une rotation autour de la verticale déplace chaque lecture et
la moyenne du même angle (\cref{exo:b1:field-lab-methods-1:11})~; la réponse est la même.

\textbf{17.} $10\times0.2 = \ang{2}$ d'erreur sur chaque \omterm{def:b1:geological-maps:attitude}{sens du pendage}, plus de
deux fois l'erreur type de \ang{0.85}~: une erreur systématique que la moyenne ne
peut pas détecter.

\textbf{18.} $80.6 - 90 = \ang{350.6}$~; l'\omterm{def:b1:geological-maps:attitude}{attitude} est direction \ang{351},
\omterm{def:b1:geological-maps:attitude}{pendage} \ang{12} vers l'est.

\textbf{19.} \omterm{def:b1:geological-maps:attitude}{Sens du pendage} $080.6$ et $080.6$~: pas de différence à cette
précision~; \omterm{def:b1:geological-maps:attitude}{pendages} $12.1$ et $12.2$~: \ang{0.1}.

\textbf{20.} Les normales $(\sin\ang{12.2}\sin\ang{80.6},
\sin\ang{12.2}\cos\ang{80.6}, \cos\ang{12.2})$ et $\vec S/\|\vec S\|$ font
un angle de \ang{0.13}.

\textbf{21.} $1.72/\sqrt{12} = \ang{0.50}$~; la différence, \ang{0.13}, en est le
quart~: non significative. La boussole et les sondages concordent.

\textbf{22.} Quatre fois plus~: 48 lectures.

\textbf{23.} Par exemple~: un réglage de déclinaison erroné ou périmé~; une
boussole déviée par du fer à proximité (un véhicule, des outils, des roches
magnétiques)~; un clinomètre déréglé~; des lectures prises seulement sur les
surfaces les mieux exposées, qui peuvent ne pas représenter le flanc. La moyenne
réduit les erreurs aléatoires, non les erreurs systématiques.

\textbf{24.} L'\omterm{def:b1:geological-maps:attitude}{attitude} moyenne du flanc est \textbf{081/12}~: \omterm{def:b1:geological-maps:attitude}{sens du pendage}
\ang{081}, \omterm{def:b1:geological-maps:attitude}{pendage} \ang{12}.
\end{solution}
